% libre-compendium
%
% Copyright (C) 2026 Richard Qin
%
% This document is dual-licensed under either:
% - GNU Free Documentation License, Version 1.3 or any later version
%   published by the Free Software Foundation; with no Invariant Sections,
%   no Front-Cover Texts, and no Back-Cover Texts (see COPYING.DOC); or
% - Creative Commons Attribution-ShareAlike 4.0 International License
%   (see LICENSE-CC-BY-SA-4.0).

% Exercises migrated from problems05.tex and solutions05.tex.
\begin{exercise}[level=basic]
设 $f(x)=x^{\mathsf T}Ax+b^{\mathsf T}x+c$，不假设 $A$ 对称。求 $\nabla f(x)$ 和 Hessian。
\begin{solution}
对标量求微分：
\begin{align*}
df
&=(dx)^{\mathsf T}Ax+x^{\mathsf T}A(dx)+b^{\mathsf T}dx\\
&=x^{\mathsf T}A^{\mathsf T}dx+x^{\mathsf T}A,dx+b^{\mathsf T}dx.
\end{align*}
因此
\[
\nabla f(x)=(A+A^{\mathsf T})x+b,
\qquad
\nabla^2f=A+A^{\mathsf T}.
\]
\end{solution}
\end{exercise}

\begin{exercise}[level=basic]
设 $f(X)=\tr(A^{\mathsf T}X)+\frac12\lVert X\rVert_F^2$。求 $\nabla_Xf$，并求使 $f$ 最小的 $X$。
\begin{solution}
\[
df=\tr(A^{\mathsf T}dX)+\tr(X^{\mathsf T}dX)
=\tr((A+X)^{\mathsf T}dX),
\]
故 $\nabla_Xf=A+X$。令梯度为零得 $X=-A$。由于二次项为
$\frac12\lVert X\rVert_F^2$，目标函数严格凸，故这是唯一最小点。
\end{solution}
\end{exercise}

\begin{exercise}[level=basic]
对 $f(x)=\frac12\lVert Ax-b\rVert_2^2$，从微分定义推导梯度和 Hessian，并说明 $A$ 列满秩时为何存在唯一最小点。
\begin{solution}
令 $r=Ax-b$，则 $dr=A,dx$，从而
\[
df=r^{\mathsf T}dr=(Ax-b)^{\mathsf T}A,dx.
\]
故
\[
\nabla f=A^{\mathsf T}(Ax-b),
\qquad
\nabla^2f=A^{\mathsf T}A.
\]
若 $A$ 列满秩，则对 $h\ne0$，
\[
h^{\mathsf T}A^{\mathsf T}Ah=\lVert Ah\rVert^2>0.
\]
因此 Hessian 正定，$f$ 严格凸，正规方程的解是唯一最小点。
\end{solution}
\end{exercise}

\begin{exercise}[level=basic]
设 $Y=AXB+C$。写出 $\operatorname{vec}(Y)$ 关于 $\operatorname{vec}(X)$ 的 Jacobian，并标明它的尺寸。
\begin{solution}
若 $A\in\R^{p\times m}$、$X\in\R^{m\times n}$、$B\in\R^{n\times q}$，则
\[
\operatorname{vec}(Y)
=(B^{\mathsf T}\otimes A)\operatorname{vec}(X)
+\operatorname{vec}(C).
\]
因此 Jacobian 为
\[
B^{\mathsf T}\otimes A\in\R^{pq\times mn}.
\]
\end{solution}
\end{exercise}

\begin{exercise}[level=basic]
设
\[
A=\begin{pmatrix}1&0\\0&0\end{pmatrix}.
\]
验证
\[
G=\begin{pmatrix}1&a\\b&c\end{pmatrix}
\]
对任意 $a,b,c\in\R$ 都是 $A$ 的一般广义逆。求其中唯一的 Moore--Penrose 逆。
\begin{solution}
直接计算
\[
AGA
=\begin{pmatrix}1&0\\0&0\end{pmatrix}
\begin{pmatrix}1&a\\b&c\end{pmatrix}
\begin{pmatrix}1&0\\0&0\end{pmatrix}
=A.
\]
因此所给 $G$ 都是一般广义逆。由 $A$ 已是奇异值对角形，其 Moore--Penrose 逆为
\[
A^+=\begin{pmatrix}1&0\\0&0\end{pmatrix}.
\]
也可由 Penrose 方程推出 $a=b=c=0$。
\end{solution}
\end{exercise}

\begin{exercise}[level=basic]
求
\[
A=\begin{pmatrix}1&1\\1&1\end{pmatrix}
\]
的 Moore--Penrose 逆，并计算 $AA^+$ 和 $A^+A$。
\begin{solution}
写成 $A=uu^{\mathsf T}$，其中 $u=(1,1)^{\mathsf T}$。秩一矩阵公式给出
\[
A^+=\frac{uu^{\mathsf T}}{\lVert u\rVert^4}
=\frac14\begin{pmatrix}1&1\\1&1\end{pmatrix}.
\]
从而
\[
AA^+=A^+A
=\frac12\begin{pmatrix}1&1\\1&1\end{pmatrix},
\]
它是到 $\operatorname{span}\{(1,1)^{\mathsf T}\}$ 的正交投影。
\end{solution}
\end{exercise}

\begin{exercise}[level=standard]
设
\[
f(X)=\frac12\lVert AXB-C\rVert_F^2
+\frac\lambda2\lVert X\rVert_F^2,
\qquad \lambda>0.
\]
求 $\nabla_Xf$，并写出驻点满足的矩阵方程。
\begin{solution}
令 $R=AXB-C$。由 $dR=A(dX)B$，
\[
df=\tr(R^{\mathsf T}A(dX)B)
+\lambda\tr(X^{\mathsf T}dX).
\]
因此
\[
\nabla_Xf=A^{\mathsf T}(AXB-C)B^{\mathsf T}+\lambda X.
\]
驻点满足
\[
A^{\mathsf T}AXB B^{\mathsf T}+\lambda X
=A^{\mathsf T}CB^{\mathsf T}.
\]
\end{solution}
\end{exercise}

\begin{exercise}[level=standard]
设 $X$ 可逆，
\[
f(X)=\log|\det X|+\tr(CX^{-1}).
\]
求 $\nabla_Xf$。
\begin{solution}
\[
d\log|\det X|=\tr(X^{-1}dX).
\]
又
\[
d\tr(CX^{-1})
=-\tr(CX^{-1}(dX)X^{-1})
=-\tr(X^{-1}CX^{-1}dX).
\]
故
\[
\nabla_Xf
=X^{-\mathsf T}-X^{-\mathsf T}C^{\mathsf T}X^{-\mathsf T}.
\]
\end{solution}
\end{exercise}

\begin{exercise}[level=standard]
设 $A,B,C$ 为给定方阵。把 Sylvester 方程
\[
AX+XB=C
\]
化为关于 $\operatorname{vec}(X)$ 的线性方程组。
\begin{solution}
利用 Vec 恒等式，
\[
\operatorname{vec}(AX)=(I\otimes A)\operatorname{vec}(X),
\]
\[
\operatorname{vec}(XB)=(B^{\mathsf T}\otimes I)\operatorname{vec}(X).
\]
因此
\[
\bigl(I\otimes A+B^{\mathsf T}\otimes I\bigr)
\operatorname{vec}(X)=\operatorname{vec}(C).
\]
\end{solution}
\end{exercise}

\begin{exercise}[level=standard]
求方程
\[
\begin{pmatrix}1&1&0\\0&1&1\end{pmatrix}x
=\begin{pmatrix}1\\1\end{pmatrix}
\]
的最小范数解。
\begin{solution}
矩阵 $A$ 行满秩，所以最小范数解为
\[
x_*=A^{\mathsf T}(AA^{\mathsf T})^{-1}b.
\]
计算
\[
AA^{\mathsf T}=\begin{pmatrix}2&1\\1&2\end{pmatrix},
\qquad
(AA^{\mathsf T})^{-1}b
=\frac13\begin{pmatrix}2&-1\\-1&2\end{pmatrix}
\begin{pmatrix}1\\1\end{pmatrix}
=\begin{pmatrix}1/3\\1/3\end{pmatrix}.
\]
因此
\[
x_*=\begin{pmatrix}1/3\\2/3\\1/3\end{pmatrix}.
\]
\end{solution}
\end{exercise}

\begin{exercise}[level=standard]
设
\[
A=\begin{pmatrix}1&0\\0&0\\0&0\end{pmatrix},
\qquad
b=\begin{pmatrix}2\\1\\-1\end{pmatrix}.
\]
求全部最小二乘解、最小范数最小二乘解和最小残差。
\begin{solution}
\[
A^+=\begin{pmatrix}1&0&0\\0&0&0\end{pmatrix},
\qquad
A^+b=\begin{pmatrix}2\\0\end{pmatrix}.
\]
又 $\ker A=\operatorname{span}\{(0,1)^{\mathsf T}\}$，故全部最小二乘解为
\[
x=\begin{pmatrix}2\\t\end{pmatrix},
\qquad t\in\R.
\]
其中最小范数解为 $(2,0)^{\mathsf T}$。投影预测为
\[
AA^+b=(2,0,0)^{\mathsf T},
\]
残差为 $(0,-1,1)^{\mathsf T}$，最小残差范数为 $\sqrt2$。
\end{solution}
\end{exercise}

\begin{exercise}[level=standard]
设 $P^2=P$。证明
\[
\R^n=\operatorname{Im}P\oplus\ker P.
\]
再证明 $P$ 是正交投影当且仅当 $P=P^{\mathsf T}$。
\begin{solution}
对任意 $x$，
\[
x=Px+(I-P)x.
\]
其中 $Px\in\operatorname{Im}P$，而
\[
P(I-P)x=(P-P^2)x=0,
\]
故 $(I-P)x\in\ker P$。若 $y\in\operatorname{Im}P\cap\ker P$，则 $y=Pz$ 且 $Py=0$，但
$Py=P^2z=Pz=y$，故 $y=0$。因此直和成立。

若 $P=P^{\mathsf T}$，对 $u=Px\in\operatorname{Im}P$、$v\in\ker P$，
\[
u^{\mathsf T}v=x^{\mathsf T}P^{\mathsf T}v=x^{\mathsf T}Pv=0,
\]
故两个子空间正交，$P$ 为正交投影。反之，正交投影在标准正交基下可写成 $Q\diag(I_r,0)Q^{\mathsf T}$，因而对称。
\end{solution}
\end{exercise}

\begin{exercise}[level=standard]
设 $A\in\R^{m\times n}$。证明
\[
\operatorname{Im}A^+=\operatorname{Im}A^{\mathsf T},
\qquad
\ker A^+=\ker A^{\mathsf T}.
\]
\begin{solution}
由 SVD $A=U\Sigma V^{\mathsf T}$，
\[
A^+=V\Sigma^+U^{\mathsf T}.
\]
$A^+$ 的像空间由 $V$ 中对应非零奇异值的右奇异向量张成，这正是 $\operatorname{Im}A^{\mathsf T}$。$A^+$ 的核由 $U$ 中对应零奇异值的向量张成，这正是
$\ker A^{\mathsf T}$。
\end{solution}
\end{exercise}

\begin{exercise}[level=standard]
设 $A$ 列满秩。证明
\[
P=A(A^{\mathsf T}A)^{-1}A^{\mathsf T}
\]
是到 $\operatorname{Im}A$ 的正交投影，并证明该正交投影与 $\operatorname{Im}A$ 的基的选择无关。
\begin{solution}
计算
\[
P^2=A(A^{\mathsf T}A)^{-1}A^{\mathsf T}A
(A^{\mathsf T}A)^{-1}A^{\mathsf T}=P,
\]
且 $P^{\mathsf T}=P$，故 $P$ 是正交投影。其像包含于 $\operatorname{Im}A$，而对 $y=Ac$，
\[
Py=A(A^{\mathsf T}A)^{-1}A^{\mathsf T}Ac=Ac=y,
\]
故像正是 $\operatorname{Im}A$。若 $B=AR$ 是同一列空间的另一组基，$R$ 可逆，则
\[
B(B^{\mathsf T}B)^{-1}B^{\mathsf T}
=AR(R^{\mathsf T}A^{\mathsf T}AR)^{-1}R^{\mathsf T}A^{\mathsf T}
=A(A^{\mathsf T}A)^{-1}A^{\mathsf T}.
\]
\end{solution}
\end{exercise}

\begin{exercise}[level=advanced]
设 $A=U\Sigma V^{\mathsf T}$ 为 SVD。不直接引用结论，验证
$A^+=V\Sigma^+U^{\mathsf T}$ 满足四个 Penrose 方程。
\begin{solution}
利用 $U^{\mathsf T}U=I$、$V^{\mathsf T}V=I$，
\[
AA^+A
=U\Sigma\Sigma^+\Sigma V^{\mathsf T}
=U\Sigma V^{\mathsf T}=A,
\]
\[
A^+AA^+
=V\Sigma^+\Sigma\Sigma^+U^{\mathsf T}=A^+.
\]
又
\[
AA^+=U(\Sigma\Sigma^+)U^{\mathsf T},
\qquad
A^+A=V(\Sigma^+\Sigma)V^{\mathsf T}.
\]
$\Sigma\Sigma^+$ 和 $\Sigma^+\Sigma$ 均为只含 $0,1$ 的对角矩阵，所以上述两矩阵均对称。四个 Penrose 方程全部成立。
\end{solution}
\end{exercise}

\begin{exercise}[level=advanced]
证明对任意 $b$，正规方程
\[
A^{\mathsf T}Ax=A^{\mathsf T}b
\]
总有解，并证明它的解集正是最小二乘解集。
\begin{solution}
$A^{\mathsf T}Ax=A^{\mathsf T}b$ 等价于
\[
A^{\mathsf T}(b-Ax)=0,
\]
即残差 $b-Ax$ 与 $\operatorname{Im}A$ 正交。正交投影定理保证存在 $y=P_{\operatorname{Im}A}b\in\operatorname{Im}A$，因此可取 $x$ 使 $Ax=y$，这证明正规方程总有解。

对任意 $x$，把残差分解为
\[
b-Ax=(b-y)+(y-Ax),
\]
其中第一项属于 $(\operatorname{Im}A)^\perp$，第二项属于 $\operatorname{Im}A$。故
\[
\lVert b-Ax\rVert^2
=\lVert b-y\rVert^2+\lVert y-Ax\rVert^2.
\]
最小值当且仅当 $Ax=y$，即当且仅当正规方程成立时取到。
\end{solution}
\end{exercise}

\begin{exercise}[level=advanced]
设 $A$ 的非零奇异值为 $\sigma_1,\ldots,\sigma_r$。证明 Tikhonov 解
\[
x_\lambda=(A^{\mathsf T}A+\lambda I)^{-1}A^{\mathsf T}b
\]
在 $\lambda\downarrow0$ 时收敛到 $A^+b$，并解释系数
$\sigma_i/(\sigma_i^2+\lambda)$ 如何抑制小奇异值方向。
\begin{solution}
由 SVD，
\[
x_\lambda
=V\diag\left(
\frac{\sigma_1}{\sigma_1^2+\lambda},\ldots,
\frac{\sigma_r}{\sigma_r^2+\lambda},0,\ldots,0
\right)U^{\mathsf T}b.
\]
对每个 $\sigma_i>0$，
\[
\frac{\sigma_i}{\sigma_i^2+\lambda}
\longrightarrow\frac1{\sigma_i},
\]
故 $x_\lambda\to V\Sigma^+U^{\mathsf T}b=A^+b$。与伪逆系数 $1/\sigma_i$ 相比，
\[
\frac{\sigma_i}{\sigma_i^2+\lambda}
=\frac1{\sigma_i}\frac{\sigma_i^2}{\sigma_i^2+\lambda}.
\]
当 $\sigma_i^2\ll\lambda$ 时，后一因子很小，因而对容易放大噪声的小奇异值方向进行衰减。
\end{solution}
\end{exercise}

\begin{exercise}[level=advanced]
给出一个满足 $P^2=P$ 但 $P\ne P^{\mathsf T}$ 的 $2\times2$ 矩阵，说明它的像空间和核空间，并解释为何它不是正交投影。
\begin{solution}
可取
\[
P=\begin{pmatrix}1&1\\0&0\end{pmatrix}.
\]
有 $P^2=P$，但 $P\ne P^{\mathsf T}$。其像空间为
$\operatorname{span}\{(1,0)^{\mathsf T}\}$，核空间由方程 $x_1+x_2=0$ 给出，即
$\operatorname{span}\{(1,-1)^{\mathsf T}\}$。这两个子空间不正交，所以 $P$ 是沿核空间方向的斜投影，不是正交投影。
\end{solution}
\end{exercise}

\begin{exercise}[level=advanced]
设
\[
A=\begin{pmatrix}1\\1\end{pmatrix},
\quad b=\begin{pmatrix}1\\3\end{pmatrix},
\quad W=\begin{pmatrix}1&0\\0&4\end{pmatrix}.
\]
求 $\min_x(Ax-b)^{\mathsf T}W(Ax-b)$ 的解，并验证残差与 $\operatorname{Im}A$ 在 $W$-内积下正交。
\begin{solution}
加权正规方程为
\[
A^{\mathsf T}WAx=A^{\mathsf T}Wb.
\]
由 $A^{\mathsf T}WA=5$、$A^{\mathsf T}Wb=13$，得
$x_*=13/5$。取残差 $r=b-Ax_*$，则
\[
r=\begin{pmatrix}-8/5\\2/5\end{pmatrix},
\qquad
A^{\mathsf T}Wr
=\begin{pmatrix}1&1\end{pmatrix}
\begin{pmatrix}1&0\\0&4\end{pmatrix}
\begin{pmatrix}-8/5\\2/5\end{pmatrix}=0.
\]
故 $r$ 与列空间在 $W$-内积下正交。
\end{solution}
\end{exercise}

\begin{exercise}[level=advanced]
用 KKT 方程求解
\[
\min_{x\in\R^2}\frac12\lVert x\rVert_2^2,
\qquad x_1+x_2=1.
\]
\begin{solution}
KKT 方程为
\[
\begin{pmatrix}1&0&1\\0&1&1\\1&1&0\end{pmatrix}
\begin{pmatrix}x_1\\x_2\\\lambda\end{pmatrix}
=\begin{pmatrix}0\\0\\1\end{pmatrix}.
\]
前两式给出 $x_1=x_2=-\lambda$，代入约束得
\[
x_1=x_2=\frac12,
\qquad \lambda=-\frac12.
\]
\end{solution}
\end{exercise}

\begin{exercise}[level=advanced]
设 $P^2=P$。证明 $P$ 的特征值只能为 $0,1$，并推出
$\rank P=\tr P$。
\begin{solution}
若 $Pv=\lambda v$ 且 $v\ne0$，则
\[
\lambda^2v=P^2v=Pv=\lambda v,
\]
故 $\lambda\in\{0,1\}$。又 $P$ 的最小多项式整除 $t(t-1)$，所以 $P$ 可对角化。对角形中 $1$ 的个数同时等于秩和迹，因此
$\rank P=\tr P$。
\end{solution}
\end{exercise}

\begin{exercise}[level=advanced]
设 $A=A^{\mathsf T}$ 的单重最大特征值为 $\lambda_1$，对应单位特征向量为 $u_1$。证明
\[
\left.\frac{d}{dt}\lambda_{\max}(A+tE)\right|_{t=0}
=u_1^{\mathsf T}Eu_1.
\]
并对
\[
A=\begin{pmatrix}2&0\\0&1\end{pmatrix},
\qquad E=\begin{pmatrix}1&2\\2&0\end{pmatrix}
\]
计算该导数。
\begin{solution}
对 $(A+tE)u(t)=\lambda(t)u(t)$ 在 $t=0$ 处求导：
\[
Eu_1+Au'(0)=\lambda'(0)u_1+\lambda_1u'(0).
\]
左乘 $u_1^{\mathsf T}$，利用
$u_1^{\mathsf T}A=\lambda_1u_1^{\mathsf T}$，得
\[
\lambda'(0)=u_1^{\mathsf T}Eu_1.
\]
题给矩阵的 $u_1=(1,0)^{\mathsf T}$，故导数为
$E_{11}=1$。
\end{solution}
\end{exercise}

\begin{exercise}[level=advanced]
考察 $A_\varepsilon=\diag(1,\varepsilon)$。求 $A_\varepsilon^+$，并说明为何 $A\mapsto A^+$ 在秩发生变化时可能不连续。
\begin{solution}
当 $\varepsilon\ne0$ 时，
\[
A_\varepsilon^+=A_\varepsilon^{-1}
=\diag(1,1/\varepsilon).
\]
当 $\varepsilon=0$ 时，$A_0^+=\diag(1,0)$。虽然
$A_\varepsilon\to A_0$，但
\[
\lVert A_\varepsilon^+\rVert_2
=\max\{1,1/|\varepsilon|\}\to\infty.
\]
因此伪逆不收敛到 $A_0^+$。极限处秩从 $2$ 降为 $1$，这正是不连续的来源。
\end{solution}
\end{exercise}

\begin{exercise}[level=advanced]
取
\[
A=\begin{pmatrix}1&1\end{pmatrix},
\qquad B=\begin{pmatrix}1\\0\end{pmatrix}.
\]
计算 $(AB)^+$ 与 $B^+A^+$，证明伪逆的逆序律一般不成立。
\begin{solution}
$AB=(1)$，故 $(AB)^+=(1)$。另一方面，
\[
A^+=\frac12\begin{pmatrix}1\\1\end{pmatrix},
\qquad B^+=\begin{pmatrix}1&0\end{pmatrix},
\]
因此
\[
B^+A^+=\frac12\ne1=(AB)^+.
\]
\end{solution}
\end{exercise}

\begin{exercise}[level=advanced]
设 $A$ 的最小非零奇异值为 $\sigma_r$。对
\[
x_*=A^+b,
\qquad \widetilde x_*=A^+(b+\delta b),
\]
证明
\[
\lVert\widetilde x_*-x_*\rVert_2
\le\frac{\lVert\delta b\rVert_2}{\sigma_r},
\]
并构造使等号成立的 $\delta b$。
\begin{solution}
\[
\widetilde x_*-x_*=A^+\delta b.
\]
由 $\lVert A^+\rVert_2=1/\sigma_r$，
\[
\lVert\widetilde x_*-x_*\rVert_2
\le\lVert A^+\rVert_2\lVert\delta b\rVert_2
=\frac{\lVert\delta b\rVert_2}{\sigma_r}.
\]
若 $u_r,v_r$ 是对应 $\sigma_r$ 的左、右奇异向量，取
$\delta b=cu_r$，则
\[
A^+\delta b=\frac{c}{\sigma_r}v_r.
\]
因此当 $c\ne0$ 时上述不等式取到等号。
\end{solution}
\end{exercise}

\begin{exercise}[level=advanced]
对
\[
A=\begin{pmatrix}1&2\\2&4\\1&2\end{pmatrix},
\]
选取一个秩分解 $A=FG$，并利用
\[
A^+=G^{\mathsf T}(GG^{\mathsf T})^{-1}
(F^{\mathsf T}F)^{-1}F^{\mathsf T}
\]
计算 $A^+$。
\begin{solution}
$A$ 的第二列是第一列的 $2$ 倍，故可取
\[
F=\begin{pmatrix}1\\2\\1\end{pmatrix},
\qquad
G=\begin{pmatrix}1&2\end{pmatrix}.
\]
此时
\[
F^{\mathsf T}F=6,
\qquad
GG^{\mathsf T}=5.
\]
因此
\begin{align*}
A^+
&=G^{\mathsf T}\frac15\frac16F^{\mathsf T}\\
&=\frac1{30}
\begin{pmatrix}1\\2\end{pmatrix}
\begin{pmatrix}1&2&1\end{pmatrix}\\
&=\frac1{30}
\begin{pmatrix}1&2&1\\2&4&2\end{pmatrix}.
\end{align*}
\end{solution}
\end{exercise}

\begin{exercise}[level=advanced]
设
\[
A(t)=\begin{pmatrix}1&t\\0&1\end{pmatrix},
\qquad
b(t)=\begin{pmatrix}t\\1\end{pmatrix},
\qquad
A(t)x(t)=b(t).
\]
不先显式求出 $A(t)^{-1}$，利用参数线性方程微分公式求 $x'(0)$。
\begin{solution}
在 $t=0$ 时，
\[
A(0)=I,
\qquad
b(0)=\begin{pmatrix}0\\1\end{pmatrix},
\qquad
x(0)=b(0).
\]
又
\[
A'(0)=\begin{pmatrix}0&1\\0&0\end{pmatrix},
\qquad
b'(0)=\begin{pmatrix}1\\0\end{pmatrix}.
\]
利用
$x'=A^{-1}(b'-A'x)$，得
\[
x'(0)
=\begin{pmatrix}1\\0\end{pmatrix}
-\begin{pmatrix}0&1\\0&0\end{pmatrix}
\begin{pmatrix}0\\1\end{pmatrix}
=\begin{pmatrix}0\\0\end{pmatrix}.
\]
\end{solution}
\end{exercise}

\begin{exercise}[level=advanced]
设
\[
A=\begin{pmatrix}1&0\\0&0\end{pmatrix},
\qquad R=A^+.
\]
令 $Y=(y_{ij})$ 为任意 $2\times2$ 矩阵。展开
\[
G=R+Y-RAYAR,
\]
并证明得到的恰是 $A$ 的全部一般广义逆。
\begin{solution}
此时 $R=A$，且 $ARA=A$、$RAR=R$。对
\[
Y=\begin{pmatrix}y_{11}&y_{12}\\y_{21}&y_{22}\end{pmatrix},
\]
有
\[
RAYAR=AYAYA
=\begin{pmatrix}y_{11}&0\\0&0\end{pmatrix}.
\]
因此
\[
G=R+Y-RAYAR
=\begin{pmatrix}1&y_{12}\\y_{21}&y_{22}\end{pmatrix}.
\]
直接计算得 $AGA=A$，所以它们都是广义逆。反之，任意满足 $AGA=A$ 的
$G=(g_{ij})$ 必满足 $g_{11}=1$，其余三个元素任意。因此上述参数化恰好覆盖全部一般广义逆。
\end{solution}
\end{exercise}

